% ============================================================
% SECTION 1 — INTRODUCTION
% ============================================================
\section{Introduction}
\label{sec:introduction}

Reliable prediction intervals are indispensable whenever
point forecasts are deployed in high-stakes settings: monetary
policy decisions hinge on the uncertainty around inflation
projections; energy-market participants require interval estimates
that correctly reflect demand volatility in winter versus summer;
supply-chain planners need intervals that are not systematically
too narrow in peak months.  In all these settings the key
feature is \emph{seasonal heteroskedasticity}: the variability
of forecast errors differs systematically by calendar month
or quarter.

The conformal-prediction literature has made substantial progress
in constructing prediction intervals for time series
\citep{xu2023conformal, xu2023sequential, barber2023conformal},
but the problem of \emph{seasonal} heteroskedasticity has
received comparatively little attention.  The standard
approach, which we call \emph{pooled calibration},
calibrates intervals on a single buffer of residuals drawn
from all seasons indiscriminately.  When the distribution of
forecast errors varies by season, pooled calibration introduces
a persistent, season-dependent bias: intervals are too narrow
(undercoverage) in volatile seasons and too wide (overcoverage)
in stable seasons, regardless of how large the training set is.

\paragraph{Contribution.}
This paper makes three contributions.

\begin{enumerate}[label=(\roman*)]
\item \emph{Formalisation of the seasonal bias.}
  We derive an explicit upper bound for the per-season
  coverage gap of the pooled EnbPI algorithm
  \citep{xu2023conformal}.  The bound separates into a
  vanishing sampling term and a non-vanishing bias term
  $2b_{s^*}$, where $b_{s^*} = \sup_x |F_{s^*}(x) - \bar{F}(x)|$
  measures the CDF discrepancy between season $s^*$ and the
  pooled residual distribution
  (Proposition~\ref{prop:pooled_gap}).

\item \emph{The \texttt{EnbPI-S} algorithm.}
  We introduce \texttt{EnbPI-S} (Algorithm~\ref{alg:enbpis}),
  a stratified extension of EnbPI that maintains $\Sper$ separate
  calibration buffers (one per calendar season).  The only
  departure from EnbPI is in the calibration step: the line
  search for the optimal asymmetric quantile pair uses only
  residuals from the same season as the prediction target.
  The training phase is unchanged; no data splitting is required;
  and the total memory footprint is identical to EnbPI.

\item \emph{Theoretical guarantees.}
  We prove that the per-season coverage gap of \texttt{EnbPI-S}
  converges to zero at rate $O(\sqrt{\log T_s / T_s})$, where
  $T_s = T/\Sper$ is the per-stratum training size
  (Theorem~\ref{thm:strat_coverage}).
  We also derive the dominance condition
  (Proposition~\ref{prop:strat_dominance})
  under which \texttt{EnbPI-S} has a strictly tighter bound
  than pooled EnbPI, and a closed-form bound on the width
  reduction achievable by stratification
  (Theorem~\ref{thm:width_bound}).
\end{enumerate}

\paragraph{Empirical illustration.}
We present two empirical applications, both using the
test period January 2015 -- December 2024.
In the first (food-at-home IPCA, BCB series 1635,
$T_s = 20$), \texttt{EnbPI-S} demonstrates clear
width adaptation (width ratio $1.78\times$ between
most and least uncertain months) consistent with
the data-limited predictions of Experiment E2.
In the second (Brazilian export growth, BCB series 1402,
$T_s = 35$), the method achieves coverage equidistant
from the nominal for both methods ($3.3$ percentage
points each), and the width ratio rises to $2.3\times$:
\texttt{EnbPI-S} assigns a $13.8\%$ interval to February
(harvest-onset, highest forecast risk) and a $6.0\%$
interval to July (mid-year, lowest forecast risk), while
the pooled method applies a constant $10.2\%$ throughout.

\paragraph{Relation to the literature.}
Conformal prediction for time series has been developed along
several orthogonal dimensions.  \citet{xu2023conformal}
introduced EnbPI, which handles non-exchangeability via a
LOO bootstrap ensemble and a sliding window update.
\citet{xu2023sequential} extended it to exploit residual
autocorrelation (SPCI), improving efficiency under serial
dependence.  \citet{barber2023conformal} established coverage
guarantees for non-exchangeable data through a weighted
conformal framework.

On the conditional validity side, Mondrian conformal
prediction \citep{vovk2005mondrian} achieves
category-conditional coverage for i.i.d.\ data by stratifying
the calibration set.  \citet{dewolf2025conditional} show that
normalised conformal prediction achieves conditional validity
for i.i.d.\ data from a location-scale family without
stratification.

\texttt{EnbPI-S} bridges these streams: it applies
Mondrian-style stratification to the non-exchangeable,
temporally dependent setting of EnbPI, providing
asymptotically valid season-conditional intervals without
requiring i.i.d.\ data, data splitting, or parametric
distributional assumptions.  Our approach is complementary
to SPCI: SPCI corrects for residual autocorrelation;
\texttt{EnbPI-S} corrects for seasonal heteroskedasticity.
The two can in principle be combined
(Section~\ref{sec:conclusion}).

The simulation study also benchmarks \texttt{EnbPI-S} against
SARIMA$(1,0,0)(1,0,0)_{12}$ with Gaussian prediction intervals,
the standard parametric alternative.
SARIMA is correctly specified for the seasonal mean of our DGP,
so its failure (September coverage of $66.2\%$ against
the $90\%$ nominal) is attributable entirely to two
compounding distributional misspecifications:
Gaussian tails (underestimating the heavy $t_3$ innovations)
and a single pooled $\hat{\sigma}$ (ignoring seasonal variance
heterogeneity).
Pooled EnbPI corrects the first failure through empirical
quantiles ($+8.7$ pp on high-volatility coverage);
\texttt{EnbPI-S} corrects both ($+5.6$ pp further),
reducing the cross-season coverage spread twelve-fold
relative to SARIMA ($0.019$ vs $0.243$).

\paragraph{Organisation.}
Section~\ref{sec:model} introduces the model and background.
Section~\ref{sec:pooled_gap} analyses the coverage gap of
pooled EnbPI.
Section~\ref{sec:algorithm} presents the \texttt{EnbPI-S}
algorithm and its properties.
Section~\ref{sec:theory} states the main theoretical results.
Section~\ref{sec:simulations} reports the simulation study.
Section~\ref{sec:empirical} presents two empirical applications:
food-at-home IPCA ($\Ts = 20$, data-limited regime) and
Brazilian export growth ($\Ts = 35$, adequate-sample regime).
Section~\ref{sec:conclusion} concludes.
All proofs are in Appendix~\ref{app:proofs}.

\paragraph{Notation.}
We use $\Sper$ for the number of seasons, $T$ for the training
size, $T_s = T/\Sper$ for the per-stratum size,
$\alpha \in (0,1)$ for the miscoverage level, and $s(t)$ for
the season of time $t$.  Empirical quantiles are denoted
$Q_p(\hat{\boldsymbol\epsilon})$ for the $p$-th quantile
of the multiset $\hat{\boldsymbol\epsilon}$.
All other notation is defined where introduced.
